\begin{wrapfigure}[3]{l}{0.1\textwidth}
	\vspace{-10pt}
	\includegraphics[width=\linewidth]{images/ikony_firefox.png}
\end{wrapfigure}

Domyślną przeglądarką internetową w systemie Ubuntu jest Firefox. Pracuje i zachowuje się ona tak samo, jak wersje na Mac OS, czy Windows. Dostępne są te same rozszerzenia, motywy i ustawienia.

\subsubsection{Przenoszenie ustawień}
Firefox przechowuje ustawienia w katalogu .mozilla, znajdującym się w katalogu domowym\footnote{Ukryte pliki i katalogi (ich nazwa zaczyna się kropką) wyświetlisz wciskając \keys{CTRL + h}.}. Wykonanie kopii zapasowej tego katalogu umożliwi przeniesienie ustawień pomiędzy różnymi systemami operacyjnymi. W innych systemach ten katalog znajduje się w podanych niżej miejscach:
\begin{itemize}
\item Windows 95, 98, ME\\
\directory{C / Windows / {Application Data} / Mozilla}\\
\directory{C / Windows/ Profiles / {Nazwa użytkownika} / {Application Data} / Mozilla}
\item Windows 2000, XP\footnote{Aby wyświetlić ukryte pliki, kliknij \menu{Narzędzia>Opcje folderów>Widok>{Pokazuj ukryte pliki i foldery}}.}\\
\directory{C / {Documents and Settings}  / {Nazwa użytkownika} / {Application Data} / Mozilla}
\item Windows Vista, 7\footnote{Patrz wyżej.}\\
\directory{C / Users / {Nazwa użytkownika} / AppData / Roaming / Mozilla}
\item Mac\\
\directory{\textasciitilde / Library / {Application Support} / Firefox}\\
\directory{\textasciitilde / Library / Mozilla / Firefox}
\end{itemize}

\subsubsection{Flash}
Wtyczka Flash została zainstalowana wraz z pakietem \textcolor{ubuntu_orange}{ubuntu-restricted-extras}. Jeżeli chcesz zainstalować samą wtyczkę Flash bez innych dodatków, potrzebujesz pakietu \textcolor{ubuntu_orange}{flashplugin-installer}. Niestety, Adobe przestało rozwijać wtyczkę Flash dla Linuksa i wersja instalowana w ten sposób jest stara. Do większości zadań wystarczy, ale wiele internetowych gier wymaga nowszej wersji. Rozwiązania są dwa. Albo instalacja przeglądarki Google Chrome (wbudowana w nią jest najnowsza wersja Flasha), albo instalacja windowsowego pluginu za pośrednictwem aplikacji Pipelight, co zostało omówione w następnej sekcji.

\subsubsection{Instalacja Pipelight}
Aby móc odtwarzać materiały przy pomocy wtyczek takich jak Silverlight, Widevine, ViewRight Caiway, czy najnowszy Flash, potrzebujesz aplikacji PipeLight. Aby ją zainstalować, wykonaj w konsoli\footnote{konsolę otworzysz kombinacją klawiszy \keys{CTRL + ALT + t}.} poniżej podane polecenia.
\begin{enumerate}
\item Dodanie repozytorium z programem PipeLight:
\begin{lstlisting}[language=bash]
sudo apt-add-repository ppa:pipelight/stable
\end{lstlisting}
\item Aktualizacja informacji o pakietach:
\begin{lstlisting}[language=bash]
sudo apt-get update
\end{lstlisting}
\item Instalacja programu:
\begin{lstlisting}[language=bash]
sudo apt-get install pipelight-multi
\end{lstlisting}
\end{enumerate}

\noindent Nie instaluj wszystkich dostępnych pluginów, gdyż znacznie spowolni to przeglądarkę. Instaluj tylko te, których potrzebujesz.

Wiele stron internetowych sprawdza zgodność przeglądarki z ich wymaganiami. Domyślnie odrzucają one nawet zgodne przeglądarki z odpowiednimi pluginami, jeżeli wykryją, że pracują na Linuksie. Aby oszukać takie strony, potrzebujesz rozszerzenia zmieniającego wartość \textcolor{ubuntu_orange}{UserAgent} wysyłaną do strony internetowej. Do Firefoksa polecany jest dodatek \href{https://addons.mozilla.org/pl-PL/firefox/addon/user-agent-overrider/}{User Agent Overrider}. Po jego zainstalowaniu otwórz \menu{{Narzędzia}>{Dodatki}>{Rozszerzenia}>{User Agent Overrider}>{Preferencje}} i wklej na końcu otwartego okna poniższe dwie linijki:
\begin{itemize}
\item Firefox 15/Windows: Mozilla/5.0 (Windows NT 6.1; WOW64; rv:15.0) Gecko/20120427 Firefox/15.0a1
\item Safari/OSX: Mozilla/5.0 (Macintosh; Intel Mac OS X 10\_7\_3) AppleWebKit/534.55.3 (KHTML, like Gecko) Version/5.1.3 Safari/534.53.10
\end{itemize}

Kiedy natrafisz na problematyczną stronę, zmień UserAgent i przeładuj ją. Przycisk pozwalający zmienić ustawienia znajduje się na prawo od paska adresu w oknie przeglądarki.

\subsubsection{Instalacja wtyczki Silverlight}
\noindent Po zainstalowaniu programu PipeLight wykonaj następujące działania:
\begin{enumerate}
\item Zainstaluj wtyczkę Silverlight:
\begin{lstlisting}[language=bash]
sudo pipelight-plugin --enable silverlight
\end{lstlisting}
\item Zostaniesz poproszony o zaakceptowanie warunków licencji. Wciśnij \keys{y}.
\item Zresetuj przeglądarkę Firefox.
\item Po ponownym uruchomieniu przeglądarki pobrane i zainstalowane zostaną potrzebne pliki. Może to potrwać kilkanaście sekund. W tym czasie przeglądarka będzie niedostępna.
\item Otwórz stronę \href{http://bubblemark.com/silverlight2.html}{testową}. Jeżeli wszystko zostało zainstalowane poprawnie, to wtyczka Silverlight wyświetli zielone, skaczące kulki.
\end{enumerate}

\subsubsection{Instalacja wtyczki Flash}
\noindent Po zainstalowaniu programu PipeLight wykonaj następujące działania:
\begin{enumerate}
\item Usuń starą wtyczkę Flash:
\begin{lstlisting}[language=bash]
sudo apt-get remove --purge flashplugin-installer
\end{lstlisting}
\item Zainstaluj nową wersję Flasha za pomocą PipeLight:
\begin{lstlisting}[language=bash]
sudo pipelight-plugin --enable flash
\end{lstlisting}
\item Zostaniesz poproszony o zaakceptowanie warunków licencji. Wciśnij \keys{y}.
\item Zresetuj przeglądarkę Firefox.
\item Po ponownym uruchomieniu przeglądarki pobrane i zainstalowane zostaną potrzebne pliki. Może to potrwać kilkanaście sekund. W tym czasie przeglądarka będzie niedostępna.
\item Otwórz stronę \href{https://www.adobe.com/software/flash/about/}{testową}. Jeżeli wszystko zostało zainstalowane poprawnie, to zobaczysz animowane logo Adobe Flash.
\end{enumerate}

\subsubsection{Instalacja wtyczki Shockwave Player}
\noindent Po zainstalowaniu programu PipeLight wykonaj następujące działania:
\begin{enumerate}
\item Zainstaluj wtyczkę:
\begin{lstlisting}[language=bash]
sudo pipelight-plugin --unlock shockwave
sudo pipelight-plugin --enable shockwave
\end{lstlisting}
\item Zostaniesz poproszony o zaakceptowanie warunków licencji. Wciśnij \keys{y}.
\item Zresetuj przeglądarkę Firefox.
\item Po ponownym uruchomieniu przeglądarki pobrane i zainstalowane zostaną potrzebne pliki. Może to potrwać kilkanaście sekund. W tym czasie przeglądarka będzie niedostępna.
\item Jeżeli na stronach wykorzystujących Shockwave Player widzisz puste prostokąty, kliknij jeden z nich prawym przyciskiem myszy i wybierz \textcolor{ubuntu_orange}{OpenGL} jako silnik renderujący.
\end{enumerate}

\subsubsection{Instalacja wtyczki Unity Webplayer}
\noindent Po zainstalowaniu programu PipeLight wykonaj następujące działania:
\begin{enumerate}
\item Zainstaluj wtyczkę:
\begin{lstlisting}[language=bash]
sudo pipelight-plugin --enable unity3d
\end{lstlisting}
\item Zostaniesz poproszony o zaakceptowanie warunków licencji. Wciśnij \keys{y}.
\item Zresetuj przeglądarkę Firefox.
\item Po ponownym uruchomieniu przeglądarki pobrane i zainstalowane zostaną potrzebne pliki. Może to potrwać nawet kilka minut. W tym czasie przeglądarka będzie niedostępna.
\item Otwórz stronę \href{https://unity3d.com/showcase/live-demos}{testową}, aby przetestować działanie wtyczki.
\end{enumerate}

\subsubsection{Instalacja wtyczki Widevine do obsługi materiałów DRM}
\noindent Po zainstalowaniu programu PipeLight wykonaj następujące działania:
\begin{enumerate}
\item Upewnij się, że masz zainstalowaną wtyczkę Flash. Wtyczka Widevine działa zarówno z wtyczką linuksową, jak i wtyczką instalowaną przez PipeLight.
\item Zainstaluj wtyczkę:
\begin{lstlisting}[language=bash]
sudo pipelight-plugin --enable widevine
\end{lstlisting}
\item Zostaniesz poproszony o zaakceptowanie warunków licencji. Wciśnij \keys{y}.
\item Zresetuj przeglądarkę Firefox.
\item Po ponownym uruchomieniu przeglądarki pobrane i zainstalowane zostaną potrzebne pliki. Może to potrwać nawet kilka minut. W tym czasie przeglądarka będzie niedostępna.
\item Otwórz stronę \href{http://www.widevine.com/demo/index.html}{testową}, aby przetestować działanie wtyczki.
\end{enumerate}

\subsubsection{Instalacja innych pluginów}
\noindent Pipelight umożliwia instalację kilku dodatkowych pluginów. Ich aktywowanie przebiega zawsze w identyczny sposób.
\begin{enumerate}
\item Zainstaluj wtyczkę:
\begin{lstlisting}[language=bash]
sudo pipelight-plugin --enable nazwa_wtyczki
\end{lstlisting}
\item Zostaniesz poproszony o zaakceptowanie warunków licencji. Wciśnij \keys{y}.
\item Zresetuj przeglądarkę Firefox.
\item Po ponownym uruchomieniu przeglądarki pobrane i zainstalowane zostaną potrzebne pliki. Może to potrwać kilkanaście sekund. W tym czasie przeglądarka będzie niedostępna.
\end{enumerate}
Lista dostępnych pluginów:
\begin{itemize}
\item \textcolor{ubuntu_orange}{adobereader} --- zastępuje wbudowaną w Firefoksa przeglądarkę PDF pluginem Adobe Acrobat Reader;
\item \textcolor{ubuntu_orange}{foxitpdf} --- zastępuje wbudowaną w Firefoksa przeglądarkę PDF pluginem Foxit Reader;
\item \textcolor{ubuntu_orange}{grandstream} --- plugin do podłączania kamer IP;
\item \textcolor{ubuntu_orange}{hikvision} --- plugin do podłączania kamer IP;
\item \textcolor{ubuntu_orange}{roblox} --- plugin do gier napisanych na silniku roblox.
\end{itemize}

\subsubsection{Rozwiązywanie problemów}
Jeżeli z jakiegokolwiek powodu przeglądarka zaczyna źle działać, to pierwszym etapem jest przywrócenie ustawień domyślnych. W ten sposób upewnisz się, czy nieprawidłowe działanie przeglądarki spowodowane jest złymi ustawieniami użytkownika (np. wadliwy dodatek), a nie kwestią systemową.
\begin{enumerate}
\item Wyłącz przeglądarkę Firefox.
\item Przenieś katalog .mozilla w bezpieczne miejsce (wciśnij \keys{CTRL + h}, aby wyświetlić ukryte pliki) lub wykonaj polecenie:
\begin{lstlisting}[language=bash]
mv ~/.mozilla ~/.mozilla_bak
\end{lstlisting}
\begin{itemize}
\item \textcolor{ubuntu_orange}{mv} --- polecenie przenoszenia plików move;
\item \textcolor{ubuntu_orange}{\textasciitilde /.mozilla} --- określenie, co ma zostać przeniesione;
	\begin{description}
	\item[\textasciitilde /] --- tylda i slash oznaczają katalog domowy użytkownika;
	\item[.mozilla] --- katalog do skopiowania;
	\end{description}
\item \textasciitilde / .mozilla.bak --- miejsce docelowe.
\end{itemize}
\item Włącz ponownie przeglądarkę i sprawdź, czy problem nadal występuje.
\end{enumerate}
